\chapter{Experimental Validation}

\section{Introduction}

The previous chapters established the theoretical basis for Bayesian Optimization (BO) with BASS models as surrogates. We discussed initialization challenges, domain discretization, and uncertainty estimation, with special attention to unexplored regions of the search space. The next step is empirical validation.

This chapter presents the experimental framework used to evaluate BASS-based BO for hyperparameter optimization. The goal is to compare BASS against established alternatives and determine where it performs well, where it underperforms, and why. In particular, we study whether BASS can address practical concerns such as cold starts, sample efficiency, and reliable optimization in high-dimensional settings.

The methods compared in this chapter are Grid Search, Random Search, Gaussian Process-based BO (GP-BO), and BASS-based BO (BASS-BO). Experiments are conducted across multiple models, datasets, and synthetic benchmark functions to provide a broad evaluation. The analysis focuses on both optimization quality and computational efficiency.

The chapter is organized as follows. First, we describe the experimental design, including methods, models, datasets, metrics, and statistical procedures. Next, we present and interpret the results. Finally, we summarize the main findings, identify limitations, and provide recommendations for future work.

\section{Experimental Design}

\subsection{Objective and Scope}

The primary objective is to evaluate BASS-BO against established hyperparameter optimization methods on two dimensions:

\begin{enumerate}
	\item \textbf{Optimization effectiveness}: the quality of the best configuration found.
	\item \textbf{Optimization efficiency}: the number of objective evaluations and runtime required to reach strong solutions.
\end{enumerate}

This evaluation criterion is consistent with the central motivation of BO and Sequential Model-Based Optimization (SMBO): reduce expensive objective function calls while maintaining strong optimization performance. The study includes both machine learning tasks and classical benchmark functions (including functions from \cite{simulationlib}) to test behavior under controlled and application-oriented conditions.

\subsection{Hyperparameter Optimization Methods}

The methods compared are:

\begin{itemize}
	\item \textbf{BASS-Based Bayesian Optimization (BASS-BO)}: the proposed method. BASS surrogates model the response surface and uncertainty, and an acquisition function selects new evaluations iteratively.
	
	\item \textbf{Grid Search}: exhaustive evaluation over a predefined grid. This method serves as a strong baseline when the search space is small, although it scales poorly.
	
	\item \textbf{Random Search}: random sampling from predefined hyperparameter distributions. It is often more efficient than Grid Search in moderate to high dimensions.
	
	\item \textbf{Gaussian Process-Based Bayesian Optimization (GP-BO)}: a standard BO baseline using Gaussian Processes for surrogate modeling and uncertainty-guided acquisition.
\end{itemize}

\subsection{Models and Datasets}

To support robust comparisons, experiments include diverse model classes and datasets.

\paragraph{Models}
\begin{itemize}
	\item \textbf{Decision Tree}: sensitive to depth, splitting criterion, and minimum sample constraints.
	\item \textbf{Support Vector Machine (SVM)}: sensitive to regularization, kernel choice, and kernel-specific parameters.
	\item \textbf{Neural Network}: sensitive to learning rate, architecture depth/width, and regularization settings.
\end{itemize}

\paragraph{Datasets}
\begin{itemize}
	\item \textbf{Synthetic datasets}: designed to test specific properties, such as nonlinearity, noise, and increasing dimensionality.
	\item \textbf{Benchmark datasets}: standard tasks from widely used repositories, including image and tabular classification/regression problems with varying complexity.
\end{itemize}

\subsection{Experimental Procedure}

Each run follows a common protocol to ensure fair comparison.

\paragraph{1) Initialization}
All methods start from the same initial design generated by random sampling. This reduces cold-start bias and ensures that early information is matched across methods.

\paragraph{2) Iterative optimization}
\begin{itemize}
	\item \textbf{BASS-BO}: fit BASS surrogate on observed trials, optimize acquisition function, evaluate selected candidate, update surrogate, repeat.
	\item \textbf{GP-BO}: same loop with GP surrogate.
	\item \textbf{Random Search}: sample and evaluate candidates independently each iteration.
	\item \textbf{Grid Search}: evaluate according to grid ordering until budget is exhausted or grid is completed.
\end{itemize}

\paragraph{3) Budget and stopping}
A fixed evaluation budget is used per task to make efficiency comparisons meaningful. When a method terminates early (for example, finite grid exhausted), remaining iterations are treated as no additional improvement.

\paragraph{4) Repetitions and randomization}
Each configuration is repeated across multiple random seeds. Method order and data split assignments are randomized to reduce systematic effects.

\subsection{Evaluation Metrics}

\paragraph{Performance metrics}
\begin{itemize}
	\item Best objective value reached within budget.
	\item Final test performance of selected configuration (for ML tasks).
	\item Area under the best-so-far curve, to capture optimization quality over time.
\end{itemize}

\paragraph{Efficiency metrics}
\begin{itemize}
	\item Number of evaluations to reach a predefined performance threshold.
	\item Wall-clock optimization time.
	\item Surrogate overhead relative to objective evaluation cost.
\end{itemize}

\subsection{Statistical Analysis}

To compare methods rigorously, we use:

\begin{itemize}
	\item \textbf{Paired comparisons} across seeds and tasks (paired t-tests when assumptions are satisfied, non-parametric alternatives otherwise).
	\item \textbf{ANOVA or mixed-effects analysis} for multi-factor effects (method, dataset type, model class).
	\item \textbf{Effect sizes} to quantify practical significance.
	\item \textbf{Confidence intervals} for key differences.
\end{itemize}

We also include robustness analyses for initialization sensitivity, discretization choices, and seed variability.

